\section{Empirical Evaluation}

\label{sec:evaluation}

\subsection{Testnet Configuration}

We deployed Nova on a 1,000-node globally distributed testnet with parameters
$k = 20$, $\alpha = 0.8$, $\beta = 11$, DAG fan-in $= 4$, and transaction arrival
rate $T = 8{,}500$ TPS.

\subsection{Results}

\begin{table}[h]
\centering
\begin{tabular}{@{}lrrrr@{}}
\toprule
\textbf{Byzantine \%} & \textbf{TPS} & \textbf{Median Latency} & \textbf{Avg Burst Size} & \textbf{P99 Latency} \\
\midrule
0\% & 8,500 & 300ms & 425 & 500ms \\
10\% & 8,200 & 400ms & 328 & 750ms \\
20\% & 7,800 & 500ms & 312 & 1,000ms \\
30\% & 6,500 & 700ms & 217 & 1,500ms \\
\bottomrule
\end{tabular}
\caption{Nova performance under varying Byzantine fractions}
\label{tab:nova-results}
\end{table}

Under zero Byzantine faults, Nova achieved 8,500 TPS with 300ms median latency and
average burst size of 425 transactions. This represents a 2.9x throughput improvement
over sequential finalization (which achieves $\lambda \cdot 1 \approx 2,900$ TPS with
the same consensus parameters).

\subsection{Burst Size Distribution}

The burst size distribution was approximately geometric with parameter $\lambda/T$,
consistent with Theorem~\ref{thm:burst-size}. The 95th percentile burst size was
2.3x the mean, indicating occasional ``mega-bursts'' when multiple frontier vertices
finalize in close succession.

%% -----------------------------------------------------------------------
%%  7. CONCLUSION
%% -----------------------------------------------------------------------
